We evaluated sampling-based semantic consistency (SMC) methods --- NLI-based and embedding-based --- on Llama-3-8B-Instruct using 1,000 balanced HaluEval QA questions. Both methods achieve chance-level AUROC ($0.4933$ and $0.4859$ respectively), a failure independently verified as not attributable to implementation error through 14 unit tests and mechanism verification.

The failure is distributional: the model produces equally consistent outputs for correct and hallucinated answers (mean SMC-NLI difference of $0.006$), indicating it operates in a systematic confabulation regime rather than the stochastic hallucination regime assumed by SMC methods. RLHF fine-tuning suppresses generative variance uniformly, eliminating the consistency-based discriminative signal.

We introduce the stochastic hallucination vs.\ systematic confabulation distinction as a necessary regime prerequisite for sampling-based detectors. This prerequisite should be empirically verified before deploying SMC methods on instruction-tuned models. When the prerequisite is violated, consistency signals are non-discriminative and alternative detection approaches --- retrieval-based verification, probing, or calibration-based methods --- should be preferred.

Our validated SMC evaluation infrastructure is released to support future work on tasks and models where the stochastic regime may hold.
